\documentclass[11pt]{article}
\usepackage[a4paper,margin=24mm]{geometry}
\usepackage{amsmath,amssymb,booktabs}
\title{Common-scale integration for bubble-weighted three-point graphs}
\author{Phi3Lab: L4E3G39, G63, G202, G394}
\date{25 October 2026}
\begin{document}
\maketitle
Use unit mass and loop measure \(d^6k/\pi^3\). Coupling and symmetry factors are excluded. The original thirty-two forests form eight groups of four bubble subtraction states. All parent projections and boundary momenta remain those of the original routed graph.

\paragraph{Two external momenta.}
With \(G=\left(\begin{smallmatrix}1&-1/2\\-1/2&1\end{smallmatrix}\right)\), the effective two-loop Gaussian matrix is a \(4\times4\) matrix. Its upper \(2\times2\) block determines \(U\), and
\[
 F=\sum_{i,j=0}^1G_{ij}\det Q_{(0,1,2+i),(0,1,2+j)}.
\]
For each three-edge subset, let \(a,b\) be the route minors with external column \(0,1\), respectively. Its positive coefficient in \(F\) is \(a^2+b^2-ab\). Thus both external directions and their cross term are retained.

\paragraph{Overall zero-order subtraction.}
For \(M=\sum\alpha_e\), \(z=F/U\), the ordinary common-scale remainder is
\[
 R(M,z)=-\log(1+z/M),\qquad
 R_M=\frac{z}{M(M+z)},\qquad
 R_{MM}=-\frac{z(2M+z)}{M^2(M+z)^2}.
\]
These signs differ from the degree-two two-point remainder. The local bubble subtraction remains degree two.

Let \(B(s)\) be the renormalized one-loop bubble, and \(A'(s)=B(s)/(1+s)^2\), \(A(0)=0\). For one weighted edge,
\[
 t_0=\frac{M-d}{d},\quad t_1=t_0+\frac zd,\quad
 J(M,d,z)=-A(t_1)+A(t_0),
\]
\[
 J_M=-\frac1d\left[\frac{B(t_1)}{(1+t_1)^2}
                        -\frac{B(t_0)}{(1+t_0)^2}\right].
\]
The initial condition is \(J(M,d,0)=0\), and
\(J_z=-B(t_1)/[d(1+t_1)^2]\).

\paragraph{Two distinct weighted edges: G63 and G202.}
Set
\[
 c_\tau=\frac{\tau-2}{(1-\tau)^2},\quad
 h_\tau=\frac1{1-\tau},\quad
 k_\tau=\frac1{(1-\tau)^2},\qquad
 \rho(\tau)=\frac{(1-4\tau)^{3/2}}6 .
\]
The complete scalar kernel is
\[
 K=c_\tau J(M,d,z)-h_\tau e J_M(M,d,z)
 +\frac{k_\tau}{\tau}J(M+e(\tau^{-1}-1),d,z),
\]
integrated with \(\rho(\tau)\,d\tau\) on the full interval \([0,1/4]\).
An independent reference retains the full nine-term product of the two rational spectral decompositions, using \(R,R_M,R_{MM}\) above.
For \(r=1-u^2\), \(\tau=(1-r^2)/4\), the positive measure is
\(u r^4\,du/6\).

\paragraph{A shared triple propagator: G39 and G394.}
Here the effective edge is \(W(s)=B(s)^2/(1+s)^3\). Its scale primitive obeys
\[
 K_z=-\frac1d W\!\left(\frac{M+z-d}{d}\right),\qquad K(0)=0,
\]
and therefore
\[
 K(M,d,z)=-d^2z\int_0^1
 \frac{B((M+vz-d)/d)^2}{(M+vz)^3}\,dv .
\]
There is no factor \(1-v\) and no \(z^2\). Those would implement a different overall Taylor order. Five exact spectral partial fractions independently establish this first-derivative identity, including the \(1/2\) coefficient of the second mass derivative for the triple propagator.

\paragraph{Measure and actual proof scope.}
The signed forest sum \(K/U^3\) is homogeneous of degree \(-6\). Sum all six sectors with one \(\alpha_e=1\) maximal and five other parameters in \([0,1]\). The normalized simplex Jacobian cancels that homogeneity; do not multiply by \(5!\). The sixth variable is the full shared-edge remainder coordinate \(v\), or the distinct-edge spectral coordinate \(u\).

The graph proof rebuilds all thirty-two forests, all 16380 nonempty edge subsets, the original Gaussian and independent positive polynomials, and both inner Schur eliminations. Scalar tests check four complete shared spectral integrals, three complete nine-term integrals, three independent mass derivatives and three complete single-edge integrals at 200-digit arithmetic with 120-digit relative agreement. These checks do not assert complete-integral numerical acceptance.

The four targets overlap the earlier cohort40. Reconcile its actual outcomes before new production; use this alternative only for targets that remain unresolved. Arithmetic interval validation, fixed diagnostics, fresh production and independent stored-request auditing remain mandatory stages.
\paragraph{Complete rational sector transport (cohort 66).}
The four graphs remain unresolved in accounting index 59. For each original
paired state, let \(\nu_P(S)\) be the smallest number of edges of \(S\) in a
monomial of \(P\). The proposal uses
\[
 \gamma(S)=\min_{\text{paired states}}
 \{ |S|+\nu_F(S)-4\nu_U(S)\},\qquad
 \beta(S)=\begin{cases}\gamma(S),&\gamma(S)\ge1,\\1/P,&\gamma(S)<1,\end{cases}
 \quad P\in\{2,4\}.
\]
These scores define a positive proposal; they are not a proof of absolute
integrability or a variance bound for the signed forest.
For every nonempty proper subset of the six effective edges, define
\[
 H(\varnothing)=1,\qquad
 H(S)=\frac{1}{\beta(S)}\sum_{e\in S}H(S\setminus\{e\}),\qquad
 Z=\sum_{e\in E}H(E\setminus\{e\}).
\]
At each step choose the next edge with probability proportional to
\(H(S\setminus\{e\})\). All \(6!=720\) orders occur, and for the order
\(\pi=(\pi_0,\ldots,\pi_5)\), with \(S_j=\{\pi_j,\ldots,\pi_5\}\),
\[
 p_\pi=\frac1Z\prod_{j=1}^5\frac1{\beta(S_j)}.
\]
Set \(\alpha_{\pi_0}=1\), \(\alpha_{\pi_j}=r_1\cdots r_j\). For each
\(j\), use
\[
 r_j=\begin{cases}
 \displaystyle\frac{c_j u_j}{1+(c_j-1)u_j},&c_j=\gamma(S_j)\ge1,\\
 u_j^P,&\gamma(S_j)<1,
 \end{cases}
 \qquad
 J_\pi=\prod_{j=1}^5\frac{dr_j}{du_j}r_j^{5-j}.
\]
If \(y\) is the independent order selector, average its six translations
\((y+k/6)\bmod1\). The complete density is
\[
 \frac16\sum_{k=0}^5\frac{J_{\pi_k}}{p_{\pi_k}}
 \rho_v(v)\sum_f\operatorname{sign}(f)
 \frac{K_f(M,d,e,F_f/U_f;v)}{U_f^3}.
\]
Here \(\rho_v=1\) for the shared-edge kernel and
\(\rho_v=v(1-v^2)^4/6\) for the distinct-edge kernel. The auxiliary
coordinate \(v\) remains unchanged: the input has one selector, five ratios
and one auxiliary coordinate, while the physical integral has six dimensions.
There is no additional factorial or primary-chart multiplicity.

The exact measure verifier rebuilds all 32 original forests and checks eight
tables, 720 orders per table, 63 conditional partitions, 41 complete-chart
moments, 35 ratio identities, and 160 boundary or selector-seam cases.
The generic eight-edge specialization also reproduces 240 previously verified
two-point orbit sectors exactly. These are actual algebraic checks; they add
no accepted complete integrals. The native implementation evaluates both
original Gaussian determinants and independently generated positive minors
over exact rationals before conversion to interval arithmetic. The independent
reference instead solves the original two-column Schur system at high
precision, retaining both external directions and their cross term.
\paragraph{One bubble and vertex forests: twenty-two further graphs.}
Inspection of the unresolved original programs finds 22 graphs with exactly
one primitive two-point bubble and only vertex-type remaining proper parts.
The new factor proof actually reconstructs 148 original forests and 74 paired
states, and exhaustively examines \(22(2^{12}-1)=90090\) nonempty edge subsets.
Every original divergent 1PI part and local Taylor order agrees.
All original-to-Schur and positive-polynomial identities pass; three
deliberate forest or Taylor-order corruptions are rejected.

Write \(\ell=\alpha_a+\alpha_b\) for the two bubble edges, and let
\(H\) be the \(5\times5\) host Gaussian matrix after a recorded affine
unimodular change of the four original loop momenta. The two independent
physical external columns remain. If \(\zeta\) is the bubble's local Taylor
variable (or one when this Taylor operator is absent), then
\[
 Q=\ell H+\alpha_a\alpha_b\zeta^2 e_1e_1^T,\qquad
 U_{\rm original}=\frac{\det Q_{\rm loop}}{\ell^2},\qquad
 F_{\rm original}=\frac1{\ell^3}
       \sum_{i,j=0}^1G_{ij}\det Q_{(0,1,2,3+i),(0,1,2,3+j)}.
\]
These equalities are checked symbolically before the bubble integration.
The two duplicated host propagators are identified in the bare graph and
kept by their original edge identities in every forest. A parent projection
that makes a different edge momentum equal does not merge that edge.

After pairing the bubble subtraction, there are three host loops and nine
effective parameters. The positive four-edge coefficients of \(F\) are
again \(a^2+b^2-ab\); \(U\) has degree three and \(F\) degree four.
With \(M=\sum_{e=1}^9\alpha_e\), \(d=\alpha_{\rm weighted}\),
and \(z_f=F_f/U_f\), the original common-scale identity gives
\[
 J(M,d,z)=A\!\left(\frac{M-d}{d}\right)
          -A\!\left(\frac{M+z-d}{d}\right),\qquad
 J_z=-\frac{B((M+z-d)/d)}{d[1+(M+z-d)/d]^2},\qquad J(M,d,0)=0.
\]
The negative derivative implements the overall three-point \(T^0\)
remainder. The bubble still has its own \(T^2\) subtraction. Retain the full
sum \(\sum_f\operatorname{sign}(f)J(M,d,z_f)/U_f^3\) before integration.

This density has degree \(-9\). The nine largest-parameter charts have
normalized-simplex Jacobian \((1+\sum_{j=1}^8x_j)^{-9}\), with no extra
\(8!\). The physical integral is eight-dimensional. The proposed exact
transport has one order selector and eight ratios, all \(9!=362880\)
orders, and nine selector translations. Replace six by nine in the earlier
conditional-mass recurrence and use
\(J_\pi=\prod_{j=1}^8r_j' r_j^{8-j}\). The score remains
\(\min_f\{|S|+\nu_{F_f}(S)-4\nu_{U_f}(S)\}\), solely as proposal design.
An independent integer enumeration checks the normalization of every order
without floating-point probabilities. The measure and subsequent numerical
stages have separate actual completion records; the successful 148-forest
factor proof does not imply acceptance of these 22 complete integrals.
\section{Closed primitives for two bubbles on a shared edge}
Define normalized primitives
\[
 A_2'(t)=\frac{B(t)^2}{(1+t)^3},\qquad
 C_2'(t)=\frac{tB(t)^2}{(1+t)^3},\qquad A_2(0)=C_2(0)=0.
\]
The complete shared-edge auxiliary integrals, with
$t_0=(M-d)/d$ and $t_1=(M+z-d)/d$, are
\[
 J_{\mathrm{E3}}=A_2(t_0)-A_2(t_1),\qquad
 J_{\mathrm{E2}}=d\{t_1[A_2(t_1)-A_2(t_0)]-[C_2(t_1)-C_2(t_0)]\}.
\]
In particular, the three-point sign is negative for $z>0$ before the
remaining forest signs are applied. These expressions integrate the
entire auxiliary coordinate; they are not approximations at a selected
external momentum.

The rationalizing substitution
\[
 w=\frac{4}{(\sqrt{t+4}+\sqrt t)^2},\qquad
 t=\frac{(1-w)^2}{w},\qquad 0<w\leq1
\]
gives
\[
 B(t)=-\frac{(1+w)^3}{6w(1-w)}\log w
       -\frac{4(w^2+w+1)}{9w}.
\]
Consequently each primitive differential is
$R_2(w)\log^2w+R_1(w)\log w+R_0(w)$ with exact rational $R_k$.
Starting at $k=2$, rational Hermite reduction writes
$R_k=H_k'+S_k$, where $S_k$ has only simple poles. Integration by
parts transfers $-kH_k/w$ into $R_{k-1}$. The final expression is
$\sum_k H_k\log^kw+\sum_k\int S_k\log^kw\,dw$.
All rational derivative reconstructions are checked exactly.

For $a\ne0$ set
\begin{align*}
 \Phi_0(w,a)&=\log(1-w/a),\\
 \Phi_1(w,a)&=\log w\log(1-w/a)+\operatorname{Li}_2(w/a),\\
 \Phi_2(w,a)&=\log^2w\log(1-w/a)
       +2\log w\operatorname{Li}_2(w/a)-2\operatorname{Li}_3(w/a).
\end{align*}
Then $\partial_w\Phi_k=\log^kw/(w-a)$. At the zero pole the primitive
is $\log^{k+1}w/(k+1)$. The non-real poles
$\rho_\pm=(1\pm i\sqrt3)/2$ are combined as conjugates, making the
answer explicitly real. Only the log-squared simple pole occurs at
$w=1$, with $\Phi_2(1,1)=-2\zeta(3)$. Subtraction of the complete
endpoint value normalizes both primitives.

For a shared-edge E3 graph this removes one physical integration
coordinate, leaving five physical projective coordinates and a separate
discrete order selector. The original child $T^2$, overall $T^0$, full
external Gram matrix and complete forest sum must remain unchanged.
Exact derivative identities and independent hypergeometric quadrature
are recorded separately from native tests and whole-integral acceptance.
The primitive derivation alone does not certify either full graph.

\section{Complete elimination of the discrete order selector}
Write $K(\alpha)$ for the complete signed shared-edge forest kernel,
including $U_f^{-3}$ and the endpoint difference above. An ordered
projective chart $\pi$ sets $\alpha_{\pi_1}=1$ and
$\alpha_{\pi_{j+1}}=r_1\cdots r_j$, $j=1,\ldots,5$. Its Jacobian is
\[
 J_\pi(u)=\prod_{j=1}^{5} r_j^{5-j}\frac{dr_j}{du_j}.
\]
If a selector chooses $\pi$ with exact probability $p_\pi$, its
importance-weighted contribution is $J_\pi K/p_\pi$. Integrating that
selector exactly gives the five-variable density
\[
 \mathcal K(u)=\sum_{\pi\in S_6}J_\pi(u)K(\alpha_\pi(u)).
\]
There is no factor $1/720$ or $1/6$, and no remaining probability
divisor. The constant chart moment is $720/5!=6$.

Let $G$ be a group of edge permutations preserving the entire signed
forest polynomial multiset, the weighted edge, and every transport
table. At the declared external Gram matrix, exact invariance permits
\[
 \mathcal K(u)=\sum_{[\pi]\in S_6/G}|[\pi]|\,
 J_\pi(u)K(\alpha_\pi(u)).
\]
Every orbit is checked for disjointness and complete coverage. This
invariance is asserted at the declared kinematics; it is not a claim
about graph automorphisms at arbitrary external momenta.

The native implementation combines the coefficients of equal rational
arguments of $A_2$ before evaluating special functions. The sum of all
argument coefficients is checked to be exactly zero, so the common
primitive normalization cancels algebraically. A separate reference
traverses all $720$ orders without quotienting, solves the original
Gaussian blocks by rational elimination, and compares every resulting
argument and coefficient exactly.

When second moments exist, complete conditional averaging obeys
$\operatorname{Var}(\mathbb E[Y\mid u])\leq\operatorname{Var}(Y)$.
This does not guarantee improvement for every finite scrambled net.
New fixed diagnostics, independent-kernel production, and stored-record
auditing are still required before accepting a complete integral.
\end{document}
